\pdfminorversion=4 % for acroread
\documentclass[t]{beamer}
%\documentclass[t,handout]{beamer}

\usepackage[english]{babel}
\usepackage[utf8]{inputenc}
\usepackage{graphicx}

\usepackage{amsmath,amssymb}
\usepackage{listings}
\lstset{frame=lines,framesep=3pt,numbers=left,numberblanklines=false,basicstyle=\ttfamily\small}

\usepackage{subfig}
\usepackage{multicol}
\usepackage{appendixnumberbeamer}

\usepackage{tikz}
\usetikzlibrary{trees} 
\usetikzlibrary{shapes.geometric}
\usetikzlibrary{positioning,shapes,shadows,arrows,calc,mindmap}
\pgfdeclarelayer{background}
\pgfdeclarelayer{foreground}
\pgfsetlayers{background,main,foreground}
\tikzstyle{activity}=[rectangle, draw=white, rounded corners, text centered, text width=8em]
\tikzstyle{data}=[rectangle, draw=white, text centered, text width=8em]
\tikzstyle{myarrow}=[->, thick, draw=white]

% Define the layers to draw the diagram
\pgfdeclarelayer{background}
\pgfdeclarelayer{foreground}
\pgfsetlayers{background,main,foreground}


\usepackage{listings}
\lstset{numbers=left,
  showstringspaces=false,
  frame={tb},
  captionpos=b,
  lineskip=0pt,
  basicstyle=\ttfamily,
  extendedchars=true,
  stepnumber=1,
  numberstyle=\small,
  xleftmargin=1em,
  breaklines
}

\renewcommand{\vec}[1]{\mathbf{#1}}
\newcommand{\lecture}{ML4AAD}
\newcommand{\lecturelong}{Machine Learning for Automated Algorithm Design}
\newcommand{\semester}{WS 2018/19}
\newcommand{\assignment}[1]{\nth{#1} Assignment}
\newcommand{\lectors}{M. Lindauer \& A. Biedenkapp}
\newcommand{\hide}[1]{}


\newcommand{\gccs}{\paragraph{General constraints for code submissions}
    
    \begin{itemize}
        \item The program can be called as stated on the exercise sheet.
        \item The program exactly returns the required output (neither less nor more) -- please use a \texttt{--verbose} option to increase the verbosity level for debugging\footnote{You might want to use \texttt{argparse} for simplicity's sake.}.
        \item Your scripts should be commented to be readable for the tutors. All functions and classes are documented with a docstring. 
        \item Provide a README ($\to$ how to install requirements and run your program(s)) and (if necessary) an installation script if your program requires any other packages.
        \item Programs are to be submitted in python $3.5$ or newer.
        \item Adding new packages to the requirements.txt is fine. If you do this however, you'll have to give a brief description why you use that package and a link to it's documentation or github page.
        \item All prepared unittests have to pass.
        \item We don't accept ipython notebook submissions.
        \item Points will be deducted if you don't fullfill these constraints.
        \item You are allowed (sometimes required) to reuse code from previous exercises.
    \end{itemize}
    \rule{\textwidth}{.5pt}
    \smallskip\\
    \noindent} 

\usetheme{Madrid}
\useinnertheme{rectangles}
\usecolortheme{whale}
\setbeamercolor{alerted text}{fg=blue}
\useoutertheme{infolines}
\setbeamertemplate{navigation symbols}{\vspace{-5pt}} % to lower the logo
\setbeamercovered{invisible}

\makeatletter
\setbeamertemplate{footline}
{
  \leavevmode%
  \hbox{%
  \begin{beamercolorbox}[wd=.333333\paperwidth,ht=2.25ex,dp=1ex,center]{author in head/foot}%
    \usebeamerfont{author in head/foot}\insertshortauthor
  \end{beamercolorbox}%
  \begin{beamercolorbox}[wd=.333333\paperwidth,ht=2.25ex,dp=1ex,center]{title in head/foot}%
    \usebeamerfont{title in head/foot}\insertshorttitle
  \end{beamercolorbox}%
  \begin{beamercolorbox}[wd=.333333\paperwidth,ht=2.25ex,dp=1ex,right]{date in head/foot}%
    \usebeamerfont{date in head/foot}\insertshortdate{}\hspace*{2em}
    \insertframenumber\hspace*{2ex} 
  \end{beamercolorbox}}%
  \vskip0pt%
}
\makeatother

\newcommand{\lit}[1]{{\footnotesize\color{black!70}[#1]}}
\newcommand{\litw}[1]{{\footnotesize\color{black!20}[#1]}}

\pgfdeclareimage[height=1.2cm]{coseallogo}{images/logos/coseal}
\pgfdeclareimage[height=1.2cm]{mlaad}{images/logos/ml4aad.png}
\pgfdeclareimage[height=1.2cm]{freiburg}{images/logos/freiburg}

\logo{\pgfuseimage{freiburg}}

\title[MLOAD]{Machine Learning and Optimization\\ for Algorithm Design}
\subtitle{ASlib}
\author{Marius Lindauer}
\institute{University of Freiburg}
\date{}
%\date{LION 2015, Lille, France}

% \AtBeginSection[] % Do nothing for \section*
% {
%   \begin{frame}{Outline}
%     \bigskip
%     \vfill
%     \tableofcontents[currentsection]
%   \end{frame}
% }


\begin{document}

% ----------------------------------------------------------------------
{
\setbeamertemplate{footline}{} % remove footer on first slide
	\frame[c]{
	\titlepage
	\begin{center}
	\vspace{-1.5cm}
	%\includegraphics[height=5em]{images/logos/freiburg}\quad\quad
	%\includegraphics[height=5em]{logo/coseal}\quad
	\includegraphics[height=5em]{images/logos/ml4aad.png}
	\end{center}
	}
}
\author{Lindauer}
\institute{}
%\logo{\pgfuseimage{mlaad}}
\logo{}

\begin{frame}[c]{How to assess the performance of algorithm selectors?}

Issues with old AS papers:

\begin{enumerate}
  \item Tedious and time-consuming task to collect data for more than $2-3$ scenarios -- only few experimental results
	\pause
  \medskip
  \item Everyone uses other scenarios -- not comparable results
  \item Runtimes are measured on different hardware -- not comparable results
  \pause
  \medskip
  \item Scenarios are not always publicly available -- not repeatable results
  \pause
  \medskip
  \item Beginners (and even experts) make mistakes -- e.g., don't consider feature costs -- invalid results
\end{enumerate}
\end{frame}
%-----------------------------------------------------------------------
%----------------------------------------------------------------------
\begin{frame}[c]{How to assess the performance of algorithm selectors?}

\begin{block}{Fair Comparison}
Algorithm selectors get the same input:
\begin{itemize}
  \item algorithm portfolio $\portfolio$
  \item training instances $\insts$
  \item instance features $\feats$
\end{itemize}
\end{block}

\begin{block}{ASlib (Algorithm Selection Library, \url{www.aslib.net})}
ASlib provides:
\begin{itemize}
  \item performance of each algorithm on each instance
  \item instance features of each instance
  \item cost to compute instance features
\end{itemize}
\end{block}

\end{frame}
%-----------------------------------------------------------------------
%----------------------------------------------------------------------
\begin{frame}[c]{ASlib $v.1.0$ Scenarios}

\centering

\small
\begin{tabular}{l ccccccc}
\hline
\hline
Scenario 				& $|I|$ &		$|A|$ 		& $\# f$  	& $\varnothing t_f$ & Ref.\\
\hline
\aspCoseal 	 			& $1294$			& $11$ 	 		& $138$ 	 & $1.3$ 	& \lit{Hoos et al. 2014} \\
\pause
\cspCoseal  	 		& $2024$			& $2$ 	 		& $17$  	& $n/a$ 		& \lit{Gent et al. 2010}	 \\
\maxsatCoseal 			& $876$				& $6$ 	 		& $37$  	& $0.1$	& \lit{Ansotegui et a. 2014} \\
\premarCoseal 			& $527$ 			& $4$ 			& $16$		 & $n/a$ & \lit{Tierney et al. 2014}\\
\proteusCoseal			& $4021$			& $22$			& $198$		 & $6.4$  & \lit{Hurley et al. 2014}\\
\qbfCoseal 	 			& $1368$			& $5$ 	 		& $46$  	 & $n/a$ 	& \lit{Kotthoff et al. 2012}\\
\sathandElevenCoseal 	& $296$ 	 		& $15$ 	 		& $115$  	 & $41.2$ 	& \lit{Xu et al. 2012}\\
\alert{\satinduElevenCoseal} 	& $300$				& $18$ 	 		& $115$  	 & $135.3$ 	& \lit{Xu et al. 2012} \\
\alert{\satrandElevenCoseal} 	& $600$ 			& $9$ 	 		& $115$  	 & $22.0$	& \lit{Xu et al. 2012} \\
\satallTwelveCoseal		& $1614$			& $31$ 	 		& $115$		 & $40.5$  & \lit{Xu et al. 2012} \\
\sathandTwelveCoseal 	& $767$				& $31$ 	 		& $115$  	 & $39.0$  & \lit{Xu et al. 2012} \\
\satinduTwelveCoseal  	& $1167$			& $31$ 	 		& $115$  	 & $80.9$  & \lit{Xu et al. 2012} \\
\satrandTwelveCoseal  	& $1362$			& $31$ 	 		& $115$ 	 & $9.0$ & \lit{Xu et al. 2012} \\
\hline
\hline
\end{tabular} 

\end{frame}
%-----------------------------------------------------------------------
%----------------------------------------------------------------------
\begin{frame}[c]{ASlib Scenarios}

\centering
\begin{tabular}{c c}
\hline
\hline
Scenario & Data Collection Time (CPU Days)\\
\hline
\aspCoseal & $25$\\
\pause
\premarCoseal  & $28$\\
\cspCoseal & $52$\\
\maxsatCoseal & $56$\\
\pause
\satinduElevenCoseal & $128$\\
\satrandElevenCoseal & $158$\\
\qbfCoseal & $163$\\
\sathandElevenCoseal & $168$\\
\pause
\sathandTwelveCoseal & $234$\\
\satinduTwelveCoseal & $284$\\
\pause
\satallTwelveCoseal & $415$\\
\satrandTwelveCoseal & $447$\\
\proteusCoseal & $596$\\
\hline
\hline
\end{tabular}

\end{frame}
%-----------------------------------------------------------------------
%----------------------------------------------------------------------
\begin{frame}[c]{ASlib Files}

\begin{description}
	\item[algorithm\_runs.arff:] performance measurements with possible repetitions and completion status of the algorithm runs
	\item[feature\_values.arff:] feature vectors for all instances
	\item[cv.arff:] splits of the instance set into training and test sets
\end{description}

(and some more\ldots)

\end{frame}
%-----------------------------------------------------------------------
%----------------------------------------------------------------------
\begin{frame}[c,fragile]{algorithm\_runs.arff}

\small
\begin{verbatim}
@RELATION ALGORITHM_RUNS_2013-SAT-Competition

@ATTRIBUTE instance_id STRING
@ATTRIBUTE repetition NUMERIC
@ATTRIBUTE algorithm STRING
@ATTRIBUTE runtime NUMERIC
@ATTRIBUTE runstatus {ok, timeout, crash, other}

@DATA
inst1.cnf, 1, lingeling, 5.0, ok
inst1.cnf, 1, clasp, 5.8, ok
inst1.cnf, 1, sparrow, 900, timeout
inst2.cnf, 1, lingeling, 588.1, ok
inst2.cnf, 1, clasp, 900,timeout
inst2.cnf, 1, sparrow, 501, ok
\end{verbatim}

\end{frame}
%-----------------------------------------------------------------------
%----------------------------------------------------------------------
\begin{frame}[c,fragile]{Reading: algorithm\_runs.arff}

\scriptsize

use \texttt{liac-arff} to read files.

\begin{verbatim}
>>> import arff
>>> data = arff.load(open("algorithm_runs.arff"))
>>> import json
>>> print(json.dumps(data))
{
  "attributes": [
    [ "instance_id", "STRING"], 
    [ "repetition", "NUMERIC"], 
    [ "algorithm", "STRING"], 
    [ "runtime", "NUMERIC"], 
    [ "runstatus", 
      ["ok","timeout","crash","other"]]
  ], 
  "relation": "ALGORITHM_RUNS_2013-SAT-Competition", 
  "description": "", 
  "data": [
    [ "inst1.cnf", 1.0, "lingeling", 5.0, "ok" ], 
    ...
    ]
}
\end{verbatim}

\end{frame}
%-----------------------------------------------------------------------
%----------------------------------------------------------------------
\begin{frame}[c,fragile]{feature\_values.arff}

\small
\begin{verbatim}
@RELATION FEATURE_VALUES_2013-SAT-Competition

@ATTRIBUTE instance_id STRING
@ATTRIBUTE repetition NUMERIC
@ATTRIBUTE number_of_variables NUMERIC
@ATTRIBUTE number_of_clauses NUMERIC
@ATTRIBUTE first_local_min_steps NUMERIC

@DATA
inst1.cnf,1,101,24,33
inst2.cnf,1,303,105,12
inst3.cnf,1,1002,337,?
\end{verbatim}

Missing values are encoded as '?'

\end{frame}
%-----------------------------------------------------------------------
%----------------------------------------------------------------------
\begin{frame}[c,fragile]{Reading: feature\_values.arff}

\scriptsize
\begin{verbatim}
>>> data = arff.load(open("feature_values.arff"))
>>> print(json.dumps(data))
{
  "attributes": [
    [ "instance_id", "STRING"], 
    [ "repetition", "NUMERIC"], 
    [ "number_of_variables", "NUMERIC"], 
    [ "number_of_clauses", "NUMERIC"], 
    [ "first_local_min_steps", "NUMERIC"]
  ], 
  "relation": "FEATURE_VALUES_2013-SAT-Competition", 
  "description": "", 
  "data": [
    [ "inst1.cnf", 1.0, 101.0, 24.0, 33.0], 
    [ "inst2.cnf", 1.0, 303.0, 105.0, 12.0], 
    [ "inst3.cnf", 1.0, 1002.0, 337.0, null]
  ]
}

\end{verbatim}

Missing values are encoded as 'None'.

\end{frame}
%-----------------------------------------------------------------------
%----------------------------------------------------------------------
\begin{frame}[c, fragile]{cv.arff}

2-fold cross validation:

\small
\begin{verbatim}
@RELATION CV_2013-SAT-Competition

@ATTRIBUTE instance_id STRING
@ATTRIBUTE repetition NUMERIC
@ATTRIBUTE fold NUMERIC

@DATA
inst1.cnf, 1, 1
inst2.cnf, 1, 1
inst3.cnf, 1, 2
inst4.cnf, 1, 2
\end{verbatim}


\end{frame}
%-----------------------------------------------------------------------
%----------------------------------------------------------------------
\begin{frame}[c,fragile]{Reading: cv.arff}

\scriptsize
\begin{verbatim}
>>> data = arff.load(open("cv.arff"))
>>> print(json.dumps(data))
{
  "attributes": [
    [ "instance_id", "STRING"], 
    [ "repetition", "NUMERIC"], 
    [ "fold", "NUMERIC"]
  ], 
  "relation": "CV_2013-SAT-Competition", 
  "description": "", 
  "data": [
    [ "inst1.cnf", 1.0, 1.0], 
    [ "inst2.cnf", 1.0, 1.0], 
    [ "inst3.cnf", 1.0, 2.0], 
    [ "inst4.cnf", 1.0, 2.0]
  ]
}


\end{verbatim}

\end{frame}
\end{document}
